% Generated by Claude Opus 5.5; no way I'm big brained enough to write these.
%
% Vertex diagram macros (\vertexDiagram, \vertexDiagramSides) for fig:VertexTable.
% Input from amspork3/amspork3_defs.tex ONLY.
% VertexTable layout: one tikz diagram per trace event constructor.
% \vertexDiagram[options]{constructor}{TlSig1}{TlSig2}{Pend1}{Pend2} (options: see "Manual overrides")
% The constructor box sits in the middle; TlSig boxes above, Pend boxes below,
% "1" boxes on the left, "2" boxes on the right. Constructor and values are math.
% An empty value means \varnothing, drawn as a minimized "label = \varnothing" box.
% Before drawing, each row (top/bottom) is measured:
%   * a box breaks after "=" only if its row would otherwise be wider than the
%     constructor box (wider box first; the other only if that makes the row fit;
%     \varnothing boxes never break);
%   * the row slides horizontally to stay within the constructor box's width
%     (centered on the constructor if it can't fit).
\makeatletter
\newsavebox\vd@tmpbox
\newdimen\vd@W \newdimen\vd@gap \vd@gap=3mm
\newdimen\vd@aOne \newdimen\vd@aTwo \newdimen\vd@bOne \newdimen\vd@bTwo
\newdimen\vd@a \newdimen\vd@b \newdimen\vd@lo \newdimen\vd@hi \newdimen\vd@tmp
\newdimen\vd@topShift \newdimen\vd@botShift
\newdimen\vd@cOne \newdimen\vd@cTwo \newdimen\vd@dOne \newdimen\vd@dTwo
\newif\ifvd@breakA \newif\ifvd@breakB
\newif\ifvd@breakTL \newif\ifvd@breakTR \newif\ifvd@breakBL \newif\ifvd@breakBR
\newcommand\vd@ifempty[3]{\if\relax\detokenize{#1}\relax #2\else #3\fi}
\newcommand\vd@fontfor[1]{\vd@ifempty{#1}{\scriptsize}{\small}}
\newcommand\vd@valfor[1]{\vd@ifempty{#1}{\varnothing}{#1}}
\tikzset{
  vertCtor/.style={draw=mainColor, text=mainColor, inner sep=3pt},
  vertVal/.style={draw=black!33, inner sep=2pt, align=left},
  vertEmpty/.style={draw=black!33, inner sep=1pt, align=left},
  vertStub/.style={draw=mainColor, line width=1pt},
}
% \vd@widths{label}{value}{one-line width}{broken width}: outer node widths
\newcommand\vd@widths[4]{%
  \sbox\vd@tmpbox{\vd@fontfor{#2}$#1 = \vd@valfor{#2}$}#3=\wd\vd@tmpbox
  \sbox\vd@tmpbox{\vd@fontfor{#2}$#1 =$}#4=\wd\vd@tmpbox
  \sbox\vd@tmpbox{\vd@fontfor{#2}$\vd@valfor{#2}$}%
  \ifdim\wd\vd@tmpbox>#4 #4=\wd\vd@tmpbox\fi
  \vd@tmp=\vd@ifempty{#2}{2.4pt}{4.4pt}% 2 * inner sep + line width
  \advance#3 by \vd@tmp \advance#4 by \vd@tmp
  \ifdim#4>#3 #4=#3\fi
  \vd@ifempty{#2}{#4=#3}{}}% \varnothing boxes never break
\newcommand\vd@tryBreakA{\ifdim\vd@aTwo<\vd@aOne \vd@a=\vd@aTwo \vd@breakAtrue\fi}
\newcommand\vd@tryBreakB{\ifdim\vd@bTwo<\vd@bOne \vd@b=\vd@bTwo \vd@breakBtrue\fi}
\newif\ifvd@tooWide
% Manual overrides, given as \vertexDiagram[options]{...} or \vertexDiagramSides[options]{...}:
%   tl break, tr break, bl break, br break = auto (default) | true | false
%   top shift, bottom shift = <dimen> (\vertexDiagram only; replaces the automatic row slide)
\def\vd@strue{true}\def\vd@sfalse{false}
\pgfkeys{/vd/.cd,
  tl break/.store in=\vd@optTL, tr break/.store in=\vd@optTR,
  bl break/.store in=\vd@optBL, br break/.store in=\vd@optBR,
  top shift/.store in=\vd@optTopShift, bottom shift/.store in=\vd@optBotShift}
\newcommand\vd@setopts[1]{%
  \def\vd@optTL{auto}\def\vd@optTR{auto}\def\vd@optBL{auto}\def\vd@optBR{auto}%
  \let\vd@optTopShift\@empty \let\vd@optBotShift\@empty
  \pgfkeys{/vd/.cd,#1}}
% \vd@override{mode macro}{if true}{if false}
\newcommand\vd@override[3]{\ifx#1\vd@strue #2\fi \ifx#1\vd@sfalse #3\fi}
\newcommand\vd@checkWidth{\vd@tmp=\vd@a \advance\vd@tmp by \vd@b \advance\vd@tmp by \vd@gap
  \ifdim\vd@tmp>\vd@W \vd@tooWidetrue \else \vd@tooWidefalse \fi}
% \vd@row{labelA}{valueA}{labelB}{valueB}{shift dimen}{modeA}{modeB}:
% sets \ifvd@breakA/B and the shift
\newcommand\vd@row[7]{%
  \vd@widths{#1}{#2}\vd@aOne\vd@aTwo \vd@widths{#3}{#4}\vd@bOne\vd@bTwo
  \vd@a=\vd@aOne \vd@b=\vd@bOne \vd@breakAfalse \vd@breakBfalse
  \vd@checkWidth
  \ifvd@tooWide
    \ifdim\vd@aOne<\vd@bOne \vd@tryBreakB \else \vd@tryBreakA \fi
    \vd@checkWidth
    \ifvd@tooWide % break the other box too, but only if that makes the row fit
      \vd@tmp=\vd@aTwo \advance\vd@tmp by \vd@bTwo \advance\vd@tmp by \vd@gap
      \ifdim\vd@tmp>\vd@W\else \vd@tryBreakA \vd@tryBreakB \fi
    \fi
  \fi
  \vd@override#6\vd@breakAtrue\vd@breakAfalse
  \vd@override#7\vd@breakBtrue\vd@breakBfalse
  \vd@a=\ifvd@breakA\vd@aTwo\else\vd@aOne\fi
  \vd@b=\ifvd@breakB\vd@bTwo\else\vd@bOne\fi
  \vd@lo=-0.5\vd@W \advance\vd@lo by 0.5\vd@gap \advance\vd@lo by \vd@a
  \vd@hi=0.5\vd@W \advance\vd@hi by -0.5\vd@gap \advance\vd@hi by -\vd@b
  \ifdim\vd@lo>\vd@hi #5=0.5\vd@lo \advance#5 by 0.5\vd@hi
  \else #5=0pt
    \ifdim#5<\vd@lo #5=\vd@lo\fi
    \ifdim#5>\vd@hi #5=\vd@hi\fi
  \fi}
% \vd@box{options}{node name}{label}{value}{break?}
\newcommand\vd@box[5]{%
  \vd@ifempty{#4}%
    {\node[vertEmpty, font=\scriptsize, #1] (#2) {$#3 = \varnothing$};}%
    {\node[vertVal, font=\small, #1] (#2) {#5{$#3 =$\\ $#4$}{$#3 = #4$}};}}
\newcommand\vd@ifTL{\ifvd@breakTL\expandafter\@firstoftwo\else\expandafter\@secondoftwo\fi}
\newcommand\vd@ifTR{\ifvd@breakTR\expandafter\@firstoftwo\else\expandafter\@secondoftwo\fi}
\newcommand\vd@ifBL{\ifvd@breakBL\expandafter\@firstoftwo\else\expandafter\@secondoftwo\fi}
\newcommand\vd@ifBR{\ifvd@breakBR\expandafter\@firstoftwo\else\expandafter\@secondoftwo\fi}
\newcommand{\vertexDiagram}[6][]{%
  \vd@setopts{#1}%
  \sbox\vd@tmpbox{$\sporktr = #2$}\vd@W=\wd\vd@tmpbox \advance\vd@W by 6.4pt
  \vd@row{\sporkTlSigA(\sporktr)}{#3}{\sporkTlSigB(\sporktr)}{#4}\vd@topShift\vd@optTL\vd@optTR
  \ifx\vd@optTopShift\@empty\else \vd@topShift=\vd@optTopShift\relax\fi
  \ifvd@breakA\vd@breakTLtrue\else\vd@breakTLfalse\fi
  \ifvd@breakB\vd@breakTRtrue\else\vd@breakTRfalse\fi
  \vd@row{\sporkPendA(\sporktr)}{#5}{\sporkPendB(\sporktr)}{#6}\vd@botShift\vd@optBL\vd@optBR
  \ifx\vd@optBotShift\@empty\else \vd@botShift=\vd@optBotShift\relax\fi
  \ifvd@breakA\vd@breakBLtrue\else\vd@breakBLfalse\fi
  \ifvd@breakB\vd@breakBRtrue\else\vd@breakBRfalse\fi
  \begin{tikzpicture}[baseline=(ctor.base)]
    \node[vertCtor] (ctor) {$\sporktr = #2$};
    \vd@box{anchor=south east, at={($(ctor.north)+(\vd@topShift-0.5\vd@gap,2.5mm)$)}}{tl}{\sporkTlSigA(\sporktr)}{#3}\vd@ifTL
    \vd@box{anchor=south west, at={($(ctor.north)+(\vd@topShift+0.5\vd@gap,2.5mm)$)}}{tr}{\sporkTlSigB(\sporktr)}{#4}\vd@ifTR
    \vd@box{anchor=north east, at={($(ctor.south)+(\vd@botShift-0.5\vd@gap,-2.5mm)$)}}{bl}{\sporkPendA(\sporktr)}{#5}\vd@ifBL
    \vd@box{anchor=north west, at={($(ctor.south)+(\vd@botShift+0.5\vd@gap,-2.5mm)$)}}{br}{\sporkPendB(\sporktr)}{#6}\vd@ifBR
    \draw[vertStub] ($(tl.south east)+(-1.5mm,0)$) -- ++(0,-2.5mm);
    \draw[vertStub] ($(tr.south west)+(1.5mm,0)$) -- ++(0,-2.5mm);
    \draw[vertStub] ($(bl.north east)+(-1.5mm,0)$) -- ++(0,2.5mm);
    \draw[vertStub] ($(br.north west)+(1.5mm,0)$) -- ++(0,2.5mm);
  \end{tikzpicture}%
}
% \vertexDiagramSides: same arguments, but the "1" quadrants sit to the left of the
% constructor and the "2" quadrants to the right (TlSig upper, Pend lower).
% Trades horizontal space for vertical. If the diagram would exceed \linewidth,
% the wider column's boxes break after "=" (then the other column if still needed).
\newcommand\vd@sideWidths{%
  \vd@a=\ifvd@breakTL\vd@aTwo\else\vd@aOne\fi
  \vd@tmp=\ifvd@breakBL\vd@bTwo\else\vd@bOne\fi
  \ifdim\vd@tmp>\vd@a \vd@a=\vd@tmp\fi
  \vd@b=\ifvd@breakTR\vd@cTwo\else\vd@cOne\fi
  \vd@tmp=\ifvd@breakBR\vd@dTwo\else\vd@dOne\fi
  \ifdim\vd@tmp>\vd@b \vd@b=\vd@tmp\fi
  \vd@tmp=\vd@W \advance\vd@tmp by \vd@a \advance\vd@tmp by \vd@b
  \advance\vd@tmp by 5mm\relax}
\newcommand{\vertexDiagramSides}[6][]{%
  \vd@setopts{#1}%
  \sbox\vd@tmpbox{$\sporktr = #2$}\vd@W=\wd\vd@tmpbox \advance\vd@W by 6.4pt
  \vd@widths{\sporkTlSigA(\sporktr)}{#3}\vd@aOne\vd@aTwo
  \vd@widths{\sporkPendA(\sporktr)}{#5}\vd@bOne\vd@bTwo
  \vd@widths{\sporkTlSigB(\sporktr)}{#4}\vd@cOne\vd@cTwo
  \vd@widths{\sporkPendB(\sporktr)}{#6}\vd@dOne\vd@dTwo
  \vd@breakTLfalse \vd@breakBLfalse \vd@breakTRfalse \vd@breakBRfalse
  \vd@sideWidths
  \ifdim\vd@tmp>\linewidth
    \ifdim\vd@a<\vd@b \vd@breakTRtrue \vd@breakBRtrue \else \vd@breakTLtrue \vd@breakBLtrue \fi
    \vd@sideWidths
    \ifdim\vd@tmp>\linewidth \vd@breakTLtrue \vd@breakBLtrue \vd@breakTRtrue \vd@breakBRtrue \fi
  \fi
  \vd@override\vd@optTL\vd@breakTLtrue\vd@breakTLfalse
  \vd@override\vd@optTR\vd@breakTRtrue\vd@breakTRfalse
  \vd@override\vd@optBL\vd@breakBLtrue\vd@breakBLfalse
  \vd@override\vd@optBR\vd@breakBRtrue\vd@breakBRfalse
  \begin{tikzpicture}[baseline=(ctor.base)]
    \node[vertCtor] (ctor) {$\sporktr = #2$};
    \vd@box{anchor=south east, at={($(ctor.west)+(-2.5mm,0.5mm)$)}}{tl}{\sporkTlSigA(\sporktr)}{#3}\vd@ifTL
    \vd@box{anchor=south west, at={($(ctor.east)+(2.5mm,0.5mm)$)}}{tr}{\sporkTlSigB(\sporktr)}{#4}\vd@ifTR
    \vd@box{anchor=north east, at={($(ctor.west)+(-2.5mm,-0.5mm)$)}}{bl}{\sporkPendA(\sporktr)}{#5}\vd@ifBL
    \vd@box{anchor=north west, at={($(ctor.east)+(2.5mm,-0.5mm)$)}}{br}{\sporkPendB(\sporktr)}{#6}\vd@ifBR
    \draw[vertStub] ($(tl.south east)+(0,1.2mm)$) -- ++(2.5mm,0);
    \draw[vertStub] ($(tr.south west)+(0,1.2mm)$) -- ++(-2.5mm,0);
    \draw[vertStub] ($(bl.north east)+(0,-1.2mm)$) -- ++(2.5mm,0);
    \draw[vertStub] ($(br.north west)+(0,-1.2mm)$) -- ++(-2.5mm,0);
  \end{tikzpicture}%
}
\makeatother
